\section{以中国人命名的曲线}
\markboth{96. 以中国人命名的曲线}{96. 以中国人命名的曲线}

前面介绍的很多曲线，常常与某个外国人的名字联系在一起。有没有用中国人的名字命名的曲线呢？当然是有的。我国的著名数学家苏步青教授，早在本世纪二十年代在大学里读书的时候，就对一些曲线的性质进行深入研究，得出了重要的结果。当时他研究过的一类曲线，在国外被叫做“Su oval”，意思就是“苏氏卵形线”。

所谓“卵形线”，就是指平面凸闭曲线。在这里，“凸”的含义可以理解成无论作这条曲线在哪一点的切线，整个曲线都在这条切线的同一侧。所以，卵形线的概念是凸多边形概念的自然推广。用 $r$ 表示卵形线在任意点 $M$ 处的曲率半径（参考第 85 节），$\varphi$ 表示 $M$ 处的切线与 $x$ 轴的夹角，可以证明对于任何卵形线都成立
\[
  \int_0^{2\pi} r \cos\varphi \dd\varphi = 0, \qquad \int_0^{2\pi} r \sin\varphi \dd\varphi = 0.
\]
如果一条卵形线满足附加条件
\[
  \int_0^{2\pi} r \cos 2\varphi \dd\varphi = 0, \qquad \int_0^{2\pi} r \sin 2\varphi \dd\varphi = 0,
\]
就称它为苏氏卵形线。

利用第 85 节计算曲率半径的公式，容易验证，椭圆、双曲线和抛物线的曲率半径都是在它们的顶点取得极值。广而言之，对于一条任意曲线 $C$，我们也把它上面使曲率半径取得极值的点叫做这条曲线的顶点。抛物线有一个顶点，双曲线有两个顶点，但它们都不是卵形线。椭圆是卵形线，并且有四个顶点。有一个著名的定理，叫做“四顶点定理”，断言任何卵形线至少有四个顶点。而对于苏氏卵形线，则可得到更好的结果：可以证明，任何苏氏卵形线至少有六个顶点。关于这个结果的证明，以及对苏氏卵形线其他性质的讨论，要用到微分几何学中一些比较专门的知识，所以这里不能逐一介绍了。
